\chapter{Asymptotyczna teoria estymacji i niepewność w praktyce}
\label{ch:asymptotic-estimation}
\paragraph{Powiązania i odsyłacze.} Wymagane prawa wielkich liczb i CTG przypomina tom I, rozdziały~\ref{I-ch:lln} i~\ref{I-ch:clt}; informację Fishera wprowadziliśmy w rozdziale~\ref{ch:fisher-information}.


\section{Dlaczego rozkład dokładny bywa nieosiągalny?}
Dla średniej z próby normalnej dysponujemy dokładnym rozkładem.
Dla MLE skomplikowanego modelu regresyjnego znamy natomiast zwykle
jedynie równanie optymalizacji. Zamiast próbować obliczać pełną
gęstość estymatora, pytamy, co dzieje się przy $n\to\infty$.
Otrzymujemy wtedy narzędzia do konstrukcji błędów standardowych,
przedziałów i testów. Nie należy jednak mylić \emph{asymptotycznej
normalności} ze zdaniem, że w dowolnej małej próbie rozkład błędu
jest w przybliżeniu normalny.

\section{Jednostajna zbieżność kryterium i zgodność estymatora}
\begin{lemat}[Jednostajne prawo wielkich liczb przy warunku Lipschitza]
\label{lem:uniform-lln}
Niech $\Theta\subset\R^d$ będzie niepustym zbiorem zwartym,
$X_1,X_2,\ldots$ niezależnymi zmiennymi o tym samym prawie $P$,
a $m:\mathsf X\times\Theta\to\R$ funkcją mierzalną względem $x$
i ciągłą względem $\theta$. Załóżmy, że dla ustalonego
$\theta_*\in\Theta$ zachodzi $\E|m(X_1,\theta_*)|<\infty$
oraz że istnieje nieujemna zmienna $L(X_1)$ z $\E L(X_1)<\infty$
taka, że dla wszystkich $\theta,\eta\in\Theta$
$|m(X_1,\theta)-m(X_1,\eta)|\le L(X_1)\|\theta-\eta\|$
prawie na pewno na jednym wspólnym zdarzeniu. Wtedy
\[
\sup_{\theta\in\Theta}\left|\frac1n\sum_{i=1}^n
m(X_i,\theta)-\E m(X_1,\theta)\right|
\longrightarrow0\quad\text{prawie na pewno}.
\]
\end{lemat}
\begin{proof}
Dzięki zwartości $\Theta$ dla każdego $\delta>0$ istnieje
skończona $\delta$-sieć $\{\theta_1,\ldots,\theta_N\}$.
Dla $\theta\in\Theta$ wybieramy węzeł odległy o mniej niż
$\delta$. Nierówność trójkąta i wspólne oszacowanie Lipschitza
dają
\[
\sup_{\theta\in\Theta}|P_nm_\theta-Pm_\theta|
\le\max_{1\le j\le N}|P_nm_{\theta_j}-Pm_{\theta_j}|
+\delta(P_nL+PL),
\]
gdzie $P_nf=n^{-1}\sum_i f(X_i)$ i $Pf=\E f(X_1)$.
Mocne prawo wielkich liczb stosowane do skończenie wielu $m_{\theta_j}$
i do $L$ zapewnia, że granica górna lewej strony jest nie większa
niż $2\delta\E L$ prawie na pewno. Bierzemy przeliczalny
ciąg $\delta=1/k$ i przecinamy odpowiadające mu zdarzenia
prawdopodobieństwa jeden; przejście $k\to\infty$ kończy dowód.
\end{proof}

\begin{twierdzenie}[Zgodność maksymalizatora]\label{thm:m-est-consistency}
Przy założeniach lematu~\ref{lem:uniform-lln} oznaczmy
$M(\theta)=\E m(X_1,\theta)$ i
$M_n(\theta)=n^{-1}\sum_i m(X_i,\theta)$.
Jeżeli $M$ ma dokładnie jeden maksymalizator $\theta_0\in\Theta$
i $\widehat\theta_n$ jest mierzalnym maksymalizatorem $M_n$
na $\Theta$ (albo mierzalnym przybliżonym maksymalizatorem
z błędem $o_P(1)$), to $\widehat\theta_n\to\theta_0$
według prawdopodobieństwa. Dla maksymalizatorów dokładnych
zachodzi nawet zbieżność prawie na pewno.
\end{twierdzenie}
\begin{proof}
Ponieważ $M$ jest ciągła na zwartej $\Theta$, dla każdego
$\varepsilon>0$ zbiór
$F_\varepsilon=\{\theta\in\Theta:\|\theta-\theta_0\|
\ge\varepsilon\}$ jest zwarty. Jeśli jest niepusty,
jedyność maksimum daje dodatnią lukę
$\Delta_\varepsilon=M(\theta_0)-\max_{F_\varepsilon}M>0$.
Na zdarzeniu $\sup_\theta|M_n-M|<\Delta_\varepsilon/3$
dla $\theta\in F_\varepsilon$ mamy
$M_n(\theta)\le M(\theta_0)-2\Delta_\varepsilon/3$,
a $M_n(\theta_0)\ge M(\theta_0)-\Delta_\varepsilon/3$.
Żaden dokładny maksymalizator nie należy więc do
$F_\varepsilon$. Lemat zapewnia tę nierówność od pewnego
indeksu prawie na pewno. Dla przybliżonych maksymalizatorów
potrzebujemy dodatkowo, by ich błąd optymalizacji był mniejszy niż
$\Delta_\varepsilon/3$; zachodzi to z prawdopodobieństwem dążącym
do jedności. Dowód jest zakończony.
\end{proof}

\begin{przyklad}[Zgodność MLE przez dywergencję KL]
W identyfikowalnym, poprawnie określonym modelu zdominowanym
$P_\theta$ średnie log-wiarygodności spełniają, przy warunkach
całkowalności, $\E_{\theta_0}\log p_\theta(X)
=\E_{\theta_0}\log p_{\theta_0}(X)
-D_{\mathrm{KL}}(P_{\theta_0}\|P_\theta)$.
Zatem funkcja graniczna ma jednoznaczne maksimum w $\theta_0$.
Aby zastosować poprzednie twierdzenie, trzeba \emph{dodatkowo}
wykazać jednostajną zbieżność po $\theta$; sama identyfikowalność
ani punktowe MPWL nie wystarczają.
\end{przyklad}

\section{Asymptotyczna normalność M-estymatorów}
\begin{twierdzenie}[Rozwinięcie równania estymacyjnego]
\label{thm:m-est-normality}
Niech $\theta_0\in\R^d$ i $\widehat\theta_n\to\theta_0$
według prawdopodobieństwa. Niech $\Psi_n:\Theta\to\R^d$
będzie losową funkcją klasy $C^1$ w ustalonym wypukłym otoczeniu
$U$ punktu $\theta_0\in U\subset\Theta$. Załóżmy, że
$\Prob(\widehat\theta_n\in U)\to1$,
$\Prob(\Psi_n(\widehat\theta_n)=0)\to1$,
$\sqrt n\Psi_n(\theta_0)\xrightarrow{d}\Normal_d(0,B)$,
gdzie $B$ jest nieujemnie określona, oraz że istnieje
ciągła funkcja $H:U\to\R^{d\times d}$ taka, że
$A:=H(\theta_0)$ jest nieosobliwa i
\[
\sup_{\theta\in U}\|D\Psi_n(\theta)-H(\theta)\|
\xrightarrow{P}0.
\]
Wtedy
\[
\sqrt n(\widehat\theta_n-\theta_0)
\xrightarrow{d}\Normal_d(0,A^{-1}B A^{-\mathsf T}).
\]
\end{twierdzenie}
\begin{proof}
Na zdarzeniu $\widehat\theta_n\in U$ i
$\Psi_n(\widehat\theta_n)=0$ wzór Newtona--Leibniza dla
funkcji wektorowej na odcinku daje
$0=\Psi_n(\theta_0)+A_n(\widehat\theta_n-\theta_0)$,
gdzie $A_n=\int_0^1D\Psi_n(\theta_0+t(\widehat\theta_n-
\theta_0))\,\dd t$.
Jednostajna zbieżność $D\Psi_n$ do $H$, ciągłość $H$ w
$\theta_0$ oraz zgodność $\widehat\theta_n$ zapewniają
$A_n\to H(\theta_0)=A$ według prawdopodobieństwa.
Z nieosobliwości $A$ wynika odwracalność
$A_n$ z prawdopodobieństwem dążącym do jedności.
Wobec tego $\sqrt n(\widehat\theta_n-\theta_0)
=-A_n^{-1}\sqrt n\Psi_n(\theta_0)$ na takim zdarzeniu.
Twierdzenie Slutsky'ego i liniowa transformacja rozkładu normalnego
dają tezę. Dowód nie wymaga założenia normalności danych.
\end{proof}

W praktyce $\Psi_n(\theta)=n^{-1}\sum_i\psi(X_i,\theta)$,
$A=\E D_\theta\psi(X_i,\theta_0)$ i
$B=\E[\psi(X_i,\theta_0)\psi(X_i,\theta_0)^{\mathsf T}]$,
przy założeniu zerowej średniej i skończonego drugiego momentu
$\psi$. Wówczas \emph{macierz kanapkowa}
$A^{-1}BA^{-\mathsf T}$ jest podstawą odpornego na błędną
specyfikację oszacowania niepewności, o ile estymand jest
właściwie określony przez równanie $\E\psi=0$.

\begin{wniosek}[Regularna asymptotyka MLE]
\label{cor:mle-asymptotic}
Niech $\Theta\subset\R^d$ będzie otwarta, a
$X_1,X_2,\ldots$ będą iid o rozkładzie $P_{\theta_0}$,
$\theta_0\in\Theta$. Niech $s_\theta(x)$ będzie score,
różniczkowalnym względem $\theta$ w wypukłym otoczeniu
$U$ punktu $\theta_0$, i załóżmy:
\begin{enumerate}
\item $\E_{\theta_0}s_{\theta_0}(X_1)=0$,
$\E_{\theta_0}\|s_{\theta_0}(X_1)\|^2<\infty$ oraz
$I(\theta_0)=\E_{\theta_0}[s_{\theta_0}s_{\theta_0}^{\mathsf T}]$
jest dodatnio określona;
\item zachodzi druga tożsamość informacyjna
$\E_{\theta_0}D_\theta s_{\theta_0}(X_1)=-I(\theta_0)$;
\item dla $\Psi_n(\theta)=n^{-1}\sum_{i=1}^n s_\theta(X_i)$
istnieje ciągła funkcja $H:U\to\R^{d\times d}$ z
$H(\theta_0)=-I(\theta_0)$, taka że
$\sup_{\theta\in U}\|D\Psi_n(\theta)-H(\theta)\|\to0$
według prawdopodobieństwa;
\item istnieje mierzalny estymator $\widehat\theta_n$ taki,
że $\widehat\theta_n\to\theta_0$ według prawdopodobieństwa,
$P(\widehat\theta_n\in U,\;\Psi_n(\widehat\theta_n)=0)\to1$.
\end{enumerate}
Wtedy
\[
\sqrt n(\widehat\theta_n-\theta_0)
\xrightarrow{d}\Normal_d(0,I(\theta_0)^{-1}).
\]
\end{wniosek}
\begin{proof}
W poprzednim twierdzeniu wybieramy $\Psi_n(\theta)=
n^{-1}\sum_i s_\theta(X_i)$. Z CLT wynika $B=I(\theta_0)$,
a z tożsamości informacyjnej
$A=\E\partial_\theta s_{\theta_0}=-I(\theta_0)$.
Po wstawieniu do macierzy kanapkowej otrzymujemy
$(-I)^{-1}I(-I)^{-\mathsf T}=I^{-1}$.
\end{proof}

\section{Metoda delty i przedziały asymptotyczne}
\begin{twierdzenie}[Wielowymiarowa metoda delty]
Niech $\sqrt n(T_n-\theta_0)\xrightarrow{d}Z$ w $\R^d$,
a $g:\R^d\to\R^k$ będzie różniczkowalna w $\theta_0$.
Wtedy
$\sqrt n(g(T_n)-g(\theta_0))\xrightarrow{d}Dg(\theta_0)Z$.
Jeśli $Z\sim\Normal_d(0,V)$, granica jest normalna o macierzy
kowariancji $Dg(\theta_0)V Dg(\theta_0)^{\mathsf T}$.
\end{twierdzenie}
\begin{proof}
Z różniczkowalności
$g(\theta_0+h)=g(\theta_0)+Dg(\theta_0)h+o(\|h\|)$.
Ciąg $\sqrt n(T_n-\theta_0)$ jest ciasny, a więc
$T_n\to\theta_0$ według prawdopodobieństwa.
Mnożąc resztę Taylora przez $\sqrt n$ dostajemy iloczyn
$\sqrt n\|T_n-\theta_0\|=O_P(1)$ i
$o_P(1)$, czyli składnik dążący do zera według
prawdopodobieństwa. Teza wynika ze Slutsky'ego.
\end{proof}

\begin{stwierdzenie}[Asymptotyczny przedział Waldowski]
\label{prop:wald-ci}
Niech $\widehat\theta_n$ estymuje rzeczywisty parametr
$\theta_0$ i załóżmy
$\sqrt n(\widehat\theta_n-\theta_0)\xrightarrow{d}
\Normal(0,V)$ dla $0<V<\infty$.
Jeśli $\widehat V_n\to V$ według prawdopodobieństwa
oraz $\widehat V_n>0$ z prawdopodobieństwem dążącym do
jedności, to przedział
$[\widehat\theta_n-z_{1-\alpha/2}\sqrt{\widehat V_n/n},
\widehat\theta_n+z_{1-\alpha/2}\sqrt{\widehat V_n/n}]$
ma asymptotyczne prawdopodobieństwo pokrycia $1-\alpha$,
gdzie $z_q$ jest kwantylem rzędu $q$ rozkładu $\Normal(0,1)$.
\end{stwierdzenie}
\begin{proof}
Z oszacowania $\widehat V_n\to V>0$ i twierdzenia Slutsky'ego
mamy $(\widehat\theta_n-\theta_0)/\sqrt{\widehat V_n/n}
\xrightarrow{d}\Normal(0,1)$. Ponieważ rozkład graniczny
jest ciągły, prawdopodobieństwo znalezienia się pomiędzy
$\pm z_{1-\alpha/2}$ dąży do $1-\alpha$; jest to dokładnie
zdarzenie pokrycia.
\end{proof}

\begin{uwaga}[Wariancja estymowana i zdarzenia wyjątkowe]
Warunek dodatniości $\widehat V_n$ z prawdopodobieństwem
dążącym do jedności nie wymaga, by oszacowanie było dodatnie
dla każdej próby. Na zdarzeniu, na którym nie jest dodatnie,
należy jawnie określić konwencję postępowania, np. nie
raportować przedziału Waldowskiego. Przy małej liczbie sukcesów
w próbie Bernoulliego metoda Wilsona zwykle zachowuje się
lepiej niż symetryczny przedział Walda.
\end{uwaga}

\begin{zadanie}[7 wad w próbie 100 sztuk: Wald czy Wilson?]
\label{zad:wilson-100}
Niech $X_i\iid\operatorname{Bernoulli}(p)$, $n=100$,
$S=7$ i $\widehat p=0{,}07$. Dla $z=1{,}96$ (poziom około
$95\%$) porównaj asymptotyczny przedział Waldowski i
przedział Wilsona uzyskany przez odwrócenie nierówności
$|\widehat p-p|\le z\sqrt{p(1-p)/n}$.
\par\emph{Rozwiązanie.} Przedział Waldowski to
$\widehat p\pm z\sqrt{\widehat p(1-\widehat p)/n}$.
Dla Wilsona przenosimy wszystkie składniki nierówności na
jedną stronę i otrzymujemy nierówność kwadratową względem $p$.
Jej pierwiastki są
\[
\frac{\widehat p+z^2/(2n)\pm
z\sqrt{\widehat p(1-\widehat p)/n+z^2/(4n^2)}}{1+z^2/n}.
\]
W odróżnieniu od Walda Wilson uwzględnia zależność wariancji
od testowanej wartości $p$ i pozostaje w $[0,1]$.
Nie jest jednak dokładnym przedziałem o nominalnym pokryciu
dla dowolnego $n$. Gdy $S=0$, Wald daje przedział zdegenerowany
w zerze, podczas gdy Wilson zachowuje dodatni górny kraniec.
\end{zadanie}


\section{Metoda delty drugiego rzędu: gdy pierwsza pochodna znika}
\begin{twierdzenie}[Jednowymiarowa metoda delty drugiego rzędu]
\label{thm:delta-second}
Niech $\sqrt n(T_n-\theta_0)\xrightarrow{d}Z$ oraz niech
$g:\R\to\R$ będzie dwukrotnie różniczkowalna w otoczeniu
$\theta_0$ z ciągłą drugą pochodną w $\theta_0$ i
$g'(\theta_0)=0$. Wtedy
\[
n\bigl(g(T_n)-g(\theta_0)\bigr)
\xrightarrow{d}\frac12g''(\theta_0)Z^2.
\]
\end{twierdzenie}
\begin{proof}
Wzór Taylora daje
$g(\theta_0+h)-g(\theta_0)=\frac12g''(\theta_0)h^2+
o(h^2)$. Ponieważ $\sqrt n(T_n-\theta_0)$ jest ciasny,
$T_n\to\theta_0$ według prawdopodobieństwa oraz
$n(T_n-\theta_0)^2=O_P(1)$. Przemnożenie rozwinięcia
przez $n$ daje resztę $o_P(1)$ i wskazaną granicę na mocy
ciągłego odwzorowania i Slutsky'ego.
\end{proof}
\begin{zadanie}[Estymacja kwadratu średniej w zerze]
Niech $X_i\iid\Normal(0,1)$ i $T_n=\bar X$. Jak zachowuje
się estymator $T_n^2$ wartości $\mu^2$ przy $\mu=0$?
\par\emph{Rozwiązanie.} Z dokładnego rozkładu
$\sqrt n\bar X\sim\Normal(0,1)$ otrzymujemy
$n\bar X^2\sim\chi_1^2$ dla każdego $n$.
Zwykła metoda delty dla $g(\mu)=\mu^2$ daje zdegenerowaną
granicę, gdyż $g'(0)=0$. Metoda drugiego rzędu odzyskuje
niegaussowski rozkład graniczny. W tym modelu warto również
rozważyć nieobciążony estymator $\bar X^2-S^2/n$,
który może przyjmować wartości ujemne: sama nieobciążoność
nie gwarantuje zachowania ograniczeń estymandu.
\end{zadanie}

\section{Model błędnie określony i wariancja kanapkowa}
\begin{zadanie}[Zliczenia z nadmierną dyspersją]
Zaobserwowano liczby zgłoszeń w dziesięciu porównywalnych
okresach: $0,0,0,0,1,1,2,2,2,10$. Oblicz estymator średniej
intensywności i porównaj błędy standardowe pod modelem Poissona
oraz bez zakładania równości średniej i wariancji.
\par\emph{Rozwiązanie.} Średnia wynosi $\bar x=1{,}8$,
a wariancja z mianownikiem $n-1$ wynosi
$s^2=9{,}0667$. Gdy przyjmiemy model
$X_i\iid\Pois(\lambda)$, MLE to $\widehat\lambda=\bar x$,
a modelowy błąd standardowy wynosi
$\sqrt{\widehat\lambda/n}\approx0{,}424$.
Jeżeli obserwacje są nadal iid, ale ich wspólny rozkład nie
musi być poissonowski, z CLT dla średniej i spójności
$s^2$ wynika oszacowanie błędu standardowego
$s/\sqrt n\approx0{,}952$.
Rozbieżność jest bardzo duża i wskazuje na nadmierną dyspersję.
Nie dowodzi sama w sobie, jaki jest jej mechanizm. Przyczyną
mogą być zmienne warunki pracy, heterogeniczność, nadmiar zer
albo obserwacja odstająca. W analizie należy porównać model
ujemnego dwumianowego, modele hierarchiczne i diagnostykę danych.
\end{zadanie}

\begin{stwierdzenie}[Kanapkowa wariancja estymatora średniej]
Niech $X_1,X_2,\ldots$ będą iid z $\E X_i=\mu$ oraz
$0<\Var(X_i)=v<\infty$. Dla równania estymacyjnego
$\Psi_n(\theta)=n^{-1}\sum_i(X_i-\theta)$ otrzymujemy
$\widehat\theta_n=\bar X$, $A=-1$, $B=v$ i
asymptotyczną wariancję $v/n$. Gdy dodatkowo błędnie
założymy model Poissona, modelowa informacja daje zamiast
tego $\mu/n$, co jest poprawne tylko przy $v=\mu$.
\end{stwierdzenie}
\begin{proof}
Score zastępczy $\psi(x,\theta)=x-\theta$ ma
$\E\psi(X,\mu)=0$, pochodną $\partial_\theta\psi=-1$,
a wariancję $\Var(\psi(X,\mu))=v$. Twierdzenie
\ref{thm:m-est-normality} daje
$\sqrt n(\bar X-\mu)\xrightarrow{d}\Normal(0,v)$.
Pod modelem Poissona $v=\mu$; poza nim równanie średniej
pozostaje sensowne, ale przy błędnej wariancji nie otrzymujemy
poprawnej kalibracji niepewności.
\end{proof}

\section{Trzy testy asymptotyczne z tej samej wiarygodności}
\begin{twierdzenie}[Jednoparametrowe twierdzenie Wilksa przy regularności]
\label{thm:wilks-one}
Niech $\theta_0$ leży we wnętrzu otwartego przedziału
$\Theta$. Dla próby iid rozważamy dwukrotnie różniczkowalną
log-wiarygodność $\ell_n$. Załóżmy istnienie zgodnego
maksymalizatora wewnętrznego $\widehat\theta_n$, takiego że
$\ell_n'(\widehat\theta_n)=0$, i niech
$\ell_n'(\theta_0)/\sqrt n\xrightarrow{d}\Normal(0,I_0)$,
gdzie $I_0>0$, oraz że istnieje wypukłe otoczenie
$U\subset\Theta$ parametru $\theta_0$ i ciągła funkcja
$J:U\to\R$ z $J(\theta_0)=I_0$, takie że
$\sup_{\theta\in U}|-\ell_n''(\theta)/n-J(\theta)|
\xrightarrow{P}0$. Wtedy
\[
2[\ell_n(\widehat\theta_n)-\ell_n(\theta_0)]
\xrightarrow{d}\chi_1^2.
\]
\end{twierdzenie}
\begin{proof}
Z twierdzenia Taylora dla pochodnej
$0=\ell_n'(\widehat\theta_n)
=\ell_n'(\theta_0)+\ell_n''(\xi_n)(\widehat\theta_n-
\theta_0)$, gdzie $\xi_n$ leży pomiędzy parametrami.
Zatem $\sqrt n(\widehat\theta_n-\theta_0)
=[\ell_n'(\theta_0)/\sqrt n]/[-\ell_n''(\xi_n)/n]$
i ma rozkład graniczny $\Normal(0,I_0^{-1})$.
Rozwijając z kolei $\ell_n(\widehat\theta_n)$ wokół
$\widehat\theta_n$, gdzie pierwsza pochodna znika, otrzymujemy
$2(\ell_n(\widehat\theta_n)-\ell_n(\theta_0))
= -\ell_n''(\zeta_n)(\widehat\theta_n-\theta_0)^2$,
przy $\zeta_n$ między $\theta_0$ a $\widehat\theta_n$.
Ponieważ $\widehat\theta_n\to\theta_0$, również punkty
$\xi_n$ i $\zeta_n$ dążą do $\theta_0$ według
prawdopodobieństwa. Ciągłość $J$ i jednostajna zbieżność
krzywizny dają zatem
$-\ell_n''(\xi_n)/n\to I_0$ i
$-\ell_n''(\zeta_n)/n\to I_0$ według
prawdopodobieństwa. Ostatnia postać ilorazu wiarygodności
zbiega więc do kwadratu zmiennej $\Normal(0,1)$,
czyli do $\chi_1^2$.
\end{proof}

\begin{przyklad}[Wald, score i LR: kiedy wyniki się rozchodzą?]
W regularnym modelu jednoparametrowym test Waldowski porównuje
$|\widehat\theta-\theta_0|$ z oszacowanym błędem standardowym,
test score używa $\ell_n'(\theta_0)$ i informacji przy hipotezie
zerowej, a test ilorazu wiarygodności używa
$2[\ell_n(\widehat\theta)-\ell_n(\theta_0)]$.
Asymptotycznie wszystkie trzy mają typową kalibrację normalną lub
$\chi^2_1$ pod $H_0$, lecz w małych próbach mogą dawać różne
wartości. Wald jest szczególnie wrażliwy na parametryzację i
przybliżenia w pobliżu brzegu. Testy dokładne i permutacyjne
omówimy w następnych częściach tomu.
\end{przyklad}

\section{Kiedy asymptotyka zawodzi? Zmienny nośnik i rzadkie zdarzenia}
\begin{twierdzenie}[Niegaussowska granica maksimum w modelu jednostajnym]
Niech $X_1,\ldots,X_n\iid U(0,\theta)$, $\theta>0$, i
$M_n=\max_iX_i$. Wtedy
$\frac n\theta(\theta-M_n)\xrightarrow{d}\operatorname{Exp}(1)$.
\end{twierdzenie}
\begin{proof}
Dla $x\ge0$ i $n>x$ mamy
\[
\Prob_\theta\left(\frac n\theta(\theta-M_n)>x\right)
=\Prob_\theta(M_n<\theta(1-x/n))=(1-x/n)^n.
\]
Granica jest $e^{-x}$, funkcją przeżycia rozkładu
wykładniczego z intensywnością jeden; w punktach $x<0$
prawdopodobieństwo przeżycia jest równe $1$.
Scharakteryzowanie funkcji rozkładu przez granice
w punktach jej ciągłości kończy dowód.
\end{proof}

Wyraz $n^{-1}$, a nie $n^{-1/2}$, jest naturalną skalą błędu
maksimum. Dlatego mechaniczne zastosowanie informacji Fishera
lub przedziału Waldowskiego byłoby błędem. Innym wyzwaniem jest
mała liczba sukcesów w modelu Bernoulliego: mimo liczebności
$n=100$ obserwacja $S=0$ lub $S=1$ może powodować bardzo słabe
działanie symetrycznego przybliżenia normalnego.

\section{Przykład obliczeniowy i tabela doboru metod}
\begin{center}\small
\begin{tabular}{>{\raggedright\arraybackslash}p{0.25\textwidth}>{\raggedright\arraybackslash}p{0.31\textwidth}>{\raggedright\arraybackslash}p{0.31\textwidth}}
\toprule
\textbf{Sytuacja} & \textbf{Pierwsza metoda} & \textbf{Weryfikacja lub alternatywa}\\
\midrule
Duża próba iid, parametr wewnętrzny & MLE i informacja Fishera & Krzywizna i diagnostyka modelu, przedział profilowy\\
Brzeg przestrzeni parametrów & Dokładna analiza rozkładu lub jednostronny przedział & Nie polegać na symetrycznym Waldzie\\
Liczba zdarzeń i ekspozycja & MLE intensywności Poissona & Dla małych liczebności przedział dokładny\\
Model może być błędnie określony & M-estymator i macierz kanapkowa & Interpretacja parametru pseudo-prawdziwego\\
Zależne obserwacje & Model zależności, korekta struktury kowariancji & Nie używać bezrefleksyjnie wariancji iid\\
Parametr zmienia nośnik & Statystyki porządkowe, rachunek dokładny & Cramér--Rao i Wilks wymagają regularności\\
\bottomrule
\end{tabular}
\end{center}

\Needspace{16\baselineskip}
\noindent\textbf{Przykład numeryczny w Pythonie (tylko biblioteka standardowa).}\par
\par
\Needspace{13\baselineskip}
\smallskip
\begin{lstlisting}[language=Python,caption={Porównanie przedziałów Walda i Wilsona dla odsetka wad.}]
from math import sqrt
n, k, z = 100, 7, 1.96
p = k / n
wald = (p - z * sqrt(p * (1-p)/n),
        p + z * sqrt(p * (1-p)/n))
den = 1 + z*z/n
center = (p + z*z/(2*n)) / den
half = z * sqrt(p*(1-p)/n + z*z/(4*n*n)) / den
wilson = (center-half, center+half)
print("Wald:", wald)
print("Wilson:", wilson)
\end{lstlisting}

\Needspace{16\baselineskip}
\noindent\textbf{Ten sam rachunek w Julii.}\par
\par
\Needspace{13\baselineskip}
\smallskip
\begin{lstlisting}[language=Julia,caption={Obliczanie przedziałów Walda i Wilsona w Julii.}]
n, k, z = 100, 7, 1.96
p = k / n
wald = (p-z*sqrt(p*(1-p)/n), p+z*sqrt(p*(1-p)/n))
den = 1 + z^2/n
center = (p + z^2/(2n)) / den
half = z*sqrt(p*(1-p)/n + z^2/(4n^2)) / den
println((wald=wald, wilson=(center-half, center+half)))
\end{lstlisting}

\paragraph{Jak używać tych wyników?}
Po pierwsze, podajemy założenie modelu niezależnych prób o wspólnym
$p$. Po drugie, raportujemy liczbę obserwacji i licznik sukcesów.
Po trzecie, sprawdzamy zachowanie przedziału na brzegu i zgodność
jego interpretacji z celem eksperymentu. Po czwarte, odróżniamy
\emph{niepewność parametru przy ustalonym modelu} od błędu
selekcji i błędów pomiaru, których żaden przedział wyznaczony z
błędnego modelu automatycznie nie naprawi.

% ================================================================
% ROZDZIAL 9 -- TESTOWANIE HIPOTEZ
% ================================================================
